%% Chapter 4, Section 4.5: Exercises

\section{Exercises}
\label{sec:4.5}

\begin{enumerate}

\item[4.1] Please download \texttt{resources.zip} from the course website.

The University of Chicago Tennis Club (UCTC) has been challenged to a doubles match by
their archrivals the Northwestern University Tennis Club (NUTC). UCTC has six players, and
you have been asked to pick the best two players to represent the club for the match.

In preparation for the match, the club players have played 100 doubles games. The file
\texttt{teams.txt} describes the different pairings in the different games. The first line
\[
  1 \quad 0 \quad 1 \quad 0 \quad -1 \quad -1
\]
indicates that players 1 and 3 (team 1) played against players 5 and 6 (team $-1$), with players 2
and 4 sitting out. The first line of the file \texttt{outcomes.txt} is $-1$, indicating that team $-1$ (players 5
and 6) won the game. You believe that each player has a skill level $(\mathrm{skill}(i) : i = 1, \ldots, 6)$. If
players $a$ and $b$ compete with players $c$ and $d$, then the probability that $a$ and $b$ win is
\[
  \psi\!\bigl(\mathrm{skill}(a) + \mathrm{skill}(b) - \mathrm{skill}(c) - \mathrm{skill}(d)\bigr);
  \qquad
  \psi(x) = \frac{1}{1 + \exp(-x)}.
\]
The function $\psi$ is called a (logistic) sigmoid function. Note that $\psi(x) + \psi(-x) = 1$; either $a$ and
$b$ win or $c$ and $d$ win. Your prior belief is that the players skills are independent and distributed
according to the normal distribution with mean 0 and variance 4. With reference to \texttt{outcomes.txt}
and \texttt{teams.txt} write down an analytical expression for the posterior distribution.

\item[4.2] In \texttt{resources.zip} there is a file \texttt{mcmc.py} with source code for a Metropolis
Hastings Markov chain.

\begin{enumerate}
\item What type of sampler is being used? Is it a good choice? Why/Why not?

\item Run the file \texttt{mcmc.py}. Does the Markov chain look like it has reached equilibrium? Explain.

\item Try modifying: the proposal distribution parameters, the length of the burn-in period, and
the number of iterations. Inspect the Markov chain output visually to assess convergence.
Give (approximate) 95\% posterior credible intervals for the skill levels.
\end{enumerate}

\item[4.3] A simple random walk on $\R^6$ is given by iteratively:
\begin{itemize}
  \item Picking $i \in \{1, 2, \ldots, 6\}$ uniformly at random;
  \item Adding a $\Normal(0,1)$ random variable to the $i$-th coordinate.
\end{itemize}
\begin{enumerate}
\item Modify the source code to use a random walk sampler. How do you expect the random
walk sampler to compare to the independence sampler in terms of efficiency?

\item Run the chain for $10^5$ iterations. Plot the sample paths and give (approximate) 95\% posterior confidence intervals for the skill levels.

\item Use Geweke's method (\texttt{geweke.py}) to assess the output from the previous part.
\end{enumerate}

\item[4.4] Consider the problem of computing the expected value of a standard Gaussian
$f = \Normal(0,1)$, which is of course 0. Let $g = \Normal(1,3)$. Using a sample size $N = 10^6$, compute the desired expected value in these three ways: (i) using rejection sampling with target $f$
and reference $g$; (ii) using importance sampling with target $f$ and proposal $g$; and (iii) using an
independence sampler with target $f$ and proposal kernel $g$. For rejection sampling and the independence sampler, compute the rejection rate. For importance sampling, compute the effective
sample size. For each of the three methods, obtain a plot similar to that in Figure~3.2.

\item[4.5] Replicate the results reported in Figure~\ref{fig:4.1} and provide plots of your results. In
addition, for the test function $h(x) = x$, estimate the integrated autocorrelation time of each of
the three chains.

\item[4.6] In this exercise, you will derive the preconditioned Crank-Nicolson (pCN)
method, a Metropolis Hastings algorithm designed to sample posterior distributions with Gaussian priors in high-dimensional settings.

\begin{enumerate}
\item Let $f(\theta) = \Normal(0, C)$ be a Gaussian prior on a parameter $\theta$. For $\beta \in (0,1)$, consider the
proposal Markov kernel $q_{\mathrm{pCN}}(\theta, \theta^*)$ implicitly defined by
\begin{equation}
  \theta^* = (1 - \beta^2)^{1/2} \theta + \beta \xi^*,
  \label{eq:4.2}
\end{equation}
where $\xi^* \sim \Normal(0, C)$. Show that $q_{\mathrm{pCN}}(\theta, \theta^*)$ satisfies detailed balance with respect to the
prior $f(\theta)$.

\item Derive and simplify the acceptance probability of a Metropolis Hastings algorithm to
sample a target posterior distribution $f(\theta \mid y) \propto f(y\mid\theta) f(\theta)$ with Gaussian prior $f(\theta) =
\Normal(0, C)$ using $q_{\mathrm{pCN}}$ as the proposal Markov kernel. Such a Metropolis Hastings algorithm
is called pCN; notice that moves from pCN are automatically accepted as long as they
increase the likelihood.
\end{enumerate}

\end{enumerate}
